% Pursuing Stacks (cotes 134-2 to 134-8): archive cote 134-7, site page 89 (PDF page 90), typescript page 531.
% typescript. Draft compared: ps6.tex:4489-4547.
% AI-transcribed by Claude Opus 5.5 (claude-opus-5-5, adaptive thinking, effort high), Message Batches API, batch ps4 (msgbatch_012Aoc1NuskATPsyBeKSYK7J), 2026-11-12; not yet checked by eye.
\scanpage{134-7/89}
\mspage{531}

do'nt allow \xout{\unsure{for|fxx}} an  amalgamated sum (or ``push-out'') construction,  except
in the trivial case when one of the two arrows to be amalgamated is an
isomorphism,  it is clear that we cannot hope to get with these models
a cofibration category.  There may be however a fibration category struc-%
ture,  and I started playing around a little\xout{x} with a possible notion of
Kan fibration -  without coming to a definite conclusion yet,  however.
I feel I should'nt dwell much longer still,  for the time being,  on get-%
ting off the ground a theory of schematic homotopy types,  maybe I'll
come back upon it later,  after next chapter,  when (hopefully) I'll be
a little more``in the game'' of Kan type conditions,  closed model category
structures  and Baues' ``halved'' variants of these.

At any rate,  whether or not a \inserted[typed]{homotopy} cofiber construction  can be carried
through for schematic homotopy types,  it seems rather clear to me that
Baues' suggestion or expectation\hand{,} about the set-up of cofibration and
fibration categories he ended up with,  is not quite justified.  Already
by the time \inserted{(in the late \unsure{sixties}) when} I first heard from Quillen about his approach towards axioma-%
tic homotopy theory,  it struck me that his approach (and the same applie\unsure{d}
to Baues ones), \xout{\unsure{is|ix}}  when applied say to semisimplicial objects in some
category A as ``models'',  is pretty strongly relying on the existence of
``enough'' projectives (or dually,  of enough injectives,  if working with
cosemisimplicial objects) in A.  When A is a topos say,  then the
category $\underline{\mathrm{Hom}}(\hand{\Delta}^{\circ},A)$ of semisimplicial objects of A does'nt carry \xout{xxx}\inserted{it seems}
(except in very special cases,  say of a totally disconnected topos) a
\inserted[typed]{reasonable}structure of a \xout{\unsure{closed}}\inserted[typed]{Quillen} model category,  nor even (I would think) of a
cofibration or a fibration category in the sense of Baues;  however I am
pretty sure that the derived category \xout{\unsure{(inxthxDx(A)}}  (with respect to the
\margin{$\times$} notion of  weak equivalence introduced by Illusie \hand{[\quad]})  has important
geometric meaning and is indeed a ``universe for homotopy theory''  - and
the first steps in developing such theory have been taken by Illusie
already in Chapter I of his thesis.  Quite similarly,  if A is an abelian
category,  Verdier's derived category  D(A) , deduced from the category
C(A) of complexes in A by localizing with respect to the set W of quasi-%
isomorphisms,  \xout{\unsure{it}} does allow for a homotopy formalism (as does any
``triangulated category''in the sense of Verdier),  however it does'nt seem
that this formalism may be deduced from a structure of a cofibration cate
-gory or a fibration category on C(A).   The same I guess holds for the
\margin{(circled) $C^{-}(A)$ \quad $D^{-}(A)$} subcategories $C^{-}(A)$ and $C^{+}(A)$,  giving rise to $D^{-}(A)$ and $D^{+}(A)$,  which
are more important still than D(A) in\xout{t} the everyday cohomology formalism
of ``spaces'' of all kinds (namely,  essentially,  of ringed topoi).  When

\margin{89}
